\begin{table}[tbp]
\centering
\small
\setlength{\tabcolsep}{4pt}
\caption{Social thresholds and the climate boundary.}
\label{tab:thresholds}
\begin{threeparttable}
\begin{tabular}{llll}
\toprule
Indicator & Unit & Threshold & Source \\
\midrule
Life expectancy & years & $\geq$ 74 & \citet{fanning2022} \\
Under-5 mortality & per 1,000 live births & $\leq$ 25 & \citet{un2015sdg} \\
Access to electricity & \% of population & $\geq$ 95 & \citet{fanning2022} \\
Employment & \% of labour force & $\geq$ 94 & \citet{fanning2022} \\
Secondary enrolment$^{a}$ & \% gross & $\geq$ 95 & \citet{fanning2022} \\
Consumption-based CO$_2$ & t per person per year & $\leq$ 1.61 & \citet{oneill2018} \\
\quad sensitivity (1.5\textdegree C) & t per person per year & $\leq$ 0.79 & \citet{ipcc2021spm} \\
\bottomrule
\end{tabular}
\begin{tablenotes}[flushleft]
\footnotesize
\item \textit{Notes:} $^{a}$Robustness check only (see Section~\ref{sec:data}). Social thresholds follow \citet{fanning2022}, except under-5 mortality, which uses SDG target 3.2. The 2\textdegree C boundary divides a 1{,}000 Gt CO$_2$ budget (2011--2100, 66\% probability) evenly over time and across 7 billion people \citep{oneill2018}. The 1.5\textdegree C value divides the IPCC remaining budget of 500 Gt CO$_2$ from 2020 (50\% probability) over 2020--2100 and the 2020 population.
\end{tablenotes}
\end{threeparttable}
\end{table}
